% fig1_esquema.tex  — Figura 1: esquema del problema VRPSPDTW-E
% Requiere en el preámbulo del documento que la incluye:
%   \usetikzlibrary{positioning,arrows.meta,calc,shapes.geometric}
% Se incluye con:  % fig1_esquema.tex  — Figura 1: esquema del problema VRPSPDTW-E
% Requiere en el preámbulo del documento que la incluye:
%   \usetikzlibrary{positioning,arrows.meta,calc,shapes.geometric}
% Se incluye con:  % fig1_esquema.tex  — Figura 1: esquema del problema VRPSPDTW-E
% Requiere en el preámbulo del documento que la incluye:
%   \usetikzlibrary{positioning,arrows.meta,calc,shapes.geometric}
% Se incluye con:  % fig1_esquema.tex  — Figura 1: esquema del problema VRPSPDTW-E
% Requiere en el preámbulo del documento que la incluye:
%   \usetikzlibrary{positioning,arrows.meta,calc,shapes.geometric}
% Se incluye con:  \input{fig1_esquema}
\begin{tikzpicture}[
  >={Stealth[length=2mm]},
  depot/.style={rectangle,draw,thick,fill=gray!20,minimum size=7mm,font=\small\bfseries},
  cust/.style={circle,draw,thick,fill=blue!10,minimum size=7mm,font=\small},
  stat/.style={regular polygon,regular polygon sides=3,draw,thick,fill=green!18,
               minimum size=8mm,inner sep=0pt,font=\scriptsize},
  arc/.style={->,very thick,gray!60!black}
]
  % ---------- Ruta (parte superior) ----------
  \node[depot] (d0) at (0,0)   {0};
  \node[cust]  (c1) at (2.2,0.7) {$c_1$};
  \node[cust]  (c2) at (4.4,1.0) {$c_2$};
  \node[stat]  (f)  at (6.6,0.3) {$f$};
  \node[cust]  (c3) at (8.8,0.8) {$c_3$};
  \node[depot] (d1) at (11,0)  {0};

  \draw[arc] (d0) -- (c1);
  \draw[arc] (c1) -- (c2);
  \draw[arc] (c2) -- (f);
  \draw[arc] (f)  -- (c3);
  \draw[arc] (c3) -- (d1);

  % etiquetas entrega (delivery) y recogida (pickup) simultaneas
  \node[font=\scriptsize,blue!50!black,above=0.5mm of c1] {$\delta_1\!\downarrow\ \rho_1\!\uparrow$};
  \node[font=\scriptsize,blue!50!black,above=0.5mm of c2] {$\delta_2\!\downarrow\ \rho_2\!\uparrow$};
  \node[font=\scriptsize,blue!50!black,above=0.5mm of c3] {$\delta_3\!\downarrow\ \rho_3\!\uparrow$};
  \node[font=\scriptsize,green!40!black,below=0.8mm of f] {recarga $q_f$};
  \node[font=\scriptsize,below=0.8mm of d0] {dep\'osito};

  % leyenda compacta
  \node[font=\scriptsize,anchor=west] at (0,2.0)
    {$\downarrow$ entrega \quad $\uparrow$ recogida \quad (simult\'aneas, una sola visita)};

  % ---------- Curva de estado de carga SoC (parte inferior) ----------
  \begin{scope}[yshift=-4.7cm]
    \draw[->] (0,0) -- (11.3,0) node[right,font=\scriptsize] {avance de la ruta};
    \draw[->] (0,0) -- (0,2.3) node[above,font=\scriptsize] {SoC $y$};
    \draw[dashed,gray] (0,2.0) -- (11,2.0) node[right,font=\scriptsize,black] {$Q$};

    % SoC: parte llena en 0, baja con la distancia, salta en la estacion f, vuelve a bajar
    \draw[very thick,green!45!black]
      (0,2.0) -- (2.2,1.45) -- (4.4,0.95) -- (6.6,0.45) % consumo hasta la estacion
      -- (6.6,1.75)                                      % salto por recarga parcial q_f
      -- (8.8,1.15) -- (11,0.55);                        % consumo hasta el deposito

    % marcas verticales de los nodos
    \foreach \x/\lab in {0/0,2.2/c_1,4.4/c_2,6.6/f,8.8/c_3,11/0}
      \draw[gray!40] (\x,0) -- (\x,-0.08) node[below,font=\scriptsize,black] {$\lab$};

    \draw[<->,green!40!black] (6.75,0.45) -- (6.75,1.75)
      node[midway,right,font=\scriptsize] {$q_f$};
    \node[font=\scriptsize,anchor=west] at (0.2,-0.7)
      {restricci\'on: $\,0 \le y \le Q$ en todo nodo (autonom\'ia)};
  \end{scope}
\end{tikzpicture}

\begin{tikzpicture}[
  >={Stealth[length=2mm]},
  depot/.style={rectangle,draw,thick,fill=gray!20,minimum size=7mm,font=\small\bfseries},
  cust/.style={circle,draw,thick,fill=blue!10,minimum size=7mm,font=\small},
  stat/.style={regular polygon,regular polygon sides=3,draw,thick,fill=green!18,
               minimum size=8mm,inner sep=0pt,font=\scriptsize},
  arc/.style={->,very thick,gray!60!black}
]
  % ---------- Ruta (parte superior) ----------
  \node[depot] (d0) at (0,0)   {0};
  \node[cust]  (c1) at (2.2,0.7) {$c_1$};
  \node[cust]  (c2) at (4.4,1.0) {$c_2$};
  \node[stat]  (f)  at (6.6,0.3) {$f$};
  \node[cust]  (c3) at (8.8,0.8) {$c_3$};
  \node[depot] (d1) at (11,0)  {0};

  \draw[arc] (d0) -- (c1);
  \draw[arc] (c1) -- (c2);
  \draw[arc] (c2) -- (f);
  \draw[arc] (f)  -- (c3);
  \draw[arc] (c3) -- (d1);

  % etiquetas entrega (delivery) y recogida (pickup) simultaneas
  \node[font=\scriptsize,blue!50!black,above=0.5mm of c1] {$\delta_1\!\downarrow\ \rho_1\!\uparrow$};
  \node[font=\scriptsize,blue!50!black,above=0.5mm of c2] {$\delta_2\!\downarrow\ \rho_2\!\uparrow$};
  \node[font=\scriptsize,blue!50!black,above=0.5mm of c3] {$\delta_3\!\downarrow\ \rho_3\!\uparrow$};
  \node[font=\scriptsize,green!40!black,below=0.8mm of f] {recarga $q_f$};
  \node[font=\scriptsize,below=0.8mm of d0] {dep\'osito};

  % leyenda compacta
  \node[font=\scriptsize,anchor=west] at (0,2.0)
    {$\downarrow$ entrega \quad $\uparrow$ recogida \quad (simult\'aneas, una sola visita)};

  % ---------- Curva de estado de carga SoC (parte inferior) ----------
  \begin{scope}[yshift=-4.7cm]
    \draw[->] (0,0) -- (11.3,0) node[right,font=\scriptsize] {avance de la ruta};
    \draw[->] (0,0) -- (0,2.3) node[above,font=\scriptsize] {SoC $y$};
    \draw[dashed,gray] (0,2.0) -- (11,2.0) node[right,font=\scriptsize,black] {$Q$};

    % SoC: parte llena en 0, baja con la distancia, salta en la estacion f, vuelve a bajar
    \draw[very thick,green!45!black]
      (0,2.0) -- (2.2,1.45) -- (4.4,0.95) -- (6.6,0.45) % consumo hasta la estacion
      -- (6.6,1.75)                                      % salto por recarga parcial q_f
      -- (8.8,1.15) -- (11,0.55);                        % consumo hasta el deposito

    % marcas verticales de los nodos
    \foreach \x/\lab in {0/0,2.2/c_1,4.4/c_2,6.6/f,8.8/c_3,11/0}
      \draw[gray!40] (\x,0) -- (\x,-0.08) node[below,font=\scriptsize,black] {$\lab$};

    \draw[<->,green!40!black] (6.75,0.45) -- (6.75,1.75)
      node[midway,right,font=\scriptsize] {$q_f$};
    \node[font=\scriptsize,anchor=west] at (0.2,-0.7)
      {restricci\'on: $\,0 \le y \le Q$ en todo nodo (autonom\'ia)};
  \end{scope}
\end{tikzpicture}

\begin{tikzpicture}[
  >={Stealth[length=2mm]},
  depot/.style={rectangle,draw,thick,fill=gray!20,minimum size=7mm,font=\small\bfseries},
  cust/.style={circle,draw,thick,fill=blue!10,minimum size=7mm,font=\small},
  stat/.style={regular polygon,regular polygon sides=3,draw,thick,fill=green!18,
               minimum size=8mm,inner sep=0pt,font=\scriptsize},
  arc/.style={->,very thick,gray!60!black}
]
  % ---------- Ruta (parte superior) ----------
  \node[depot] (d0) at (0,0)   {0};
  \node[cust]  (c1) at (2.2,0.7) {$c_1$};
  \node[cust]  (c2) at (4.4,1.0) {$c_2$};
  \node[stat]  (f)  at (6.6,0.3) {$f$};
  \node[cust]  (c3) at (8.8,0.8) {$c_3$};
  \node[depot] (d1) at (11,0)  {0};

  \draw[arc] (d0) -- (c1);
  \draw[arc] (c1) -- (c2);
  \draw[arc] (c2) -- (f);
  \draw[arc] (f)  -- (c3);
  \draw[arc] (c3) -- (d1);

  % etiquetas entrega (delivery) y recogida (pickup) simultaneas
  \node[font=\scriptsize,blue!50!black,above=0.5mm of c1] {$\delta_1\!\downarrow\ \rho_1\!\uparrow$};
  \node[font=\scriptsize,blue!50!black,above=0.5mm of c2] {$\delta_2\!\downarrow\ \rho_2\!\uparrow$};
  \node[font=\scriptsize,blue!50!black,above=0.5mm of c3] {$\delta_3\!\downarrow\ \rho_3\!\uparrow$};
  \node[font=\scriptsize,green!40!black,below=0.8mm of f] {recarga $q_f$};
  \node[font=\scriptsize,below=0.8mm of d0] {dep\'osito};

  % leyenda compacta
  \node[font=\scriptsize,anchor=west] at (0,2.0)
    {$\downarrow$ entrega \quad $\uparrow$ recogida \quad (simult\'aneas, una sola visita)};

  % ---------- Curva de estado de carga SoC (parte inferior) ----------
  \begin{scope}[yshift=-4.7cm]
    \draw[->] (0,0) -- (11.3,0) node[right,font=\scriptsize] {avance de la ruta};
    \draw[->] (0,0) -- (0,2.3) node[above,font=\scriptsize] {SoC $y$};
    \draw[dashed,gray] (0,2.0) -- (11,2.0) node[right,font=\scriptsize,black] {$Q$};

    % SoC: parte llena en 0, baja con la distancia, salta en la estacion f, vuelve a bajar
    \draw[very thick,green!45!black]
      (0,2.0) -- (2.2,1.45) -- (4.4,0.95) -- (6.6,0.45) % consumo hasta la estacion
      -- (6.6,1.75)                                      % salto por recarga parcial q_f
      -- (8.8,1.15) -- (11,0.55);                        % consumo hasta el deposito

    % marcas verticales de los nodos
    \foreach \x/\lab in {0/0,2.2/c_1,4.4/c_2,6.6/f,8.8/c_3,11/0}
      \draw[gray!40] (\x,0) -- (\x,-0.08) node[below,font=\scriptsize,black] {$\lab$};

    \draw[<->,green!40!black] (6.75,0.45) -- (6.75,1.75)
      node[midway,right,font=\scriptsize] {$q_f$};
    \node[font=\scriptsize,anchor=west] at (0.2,-0.7)
      {restricci\'on: $\,0 \le y \le Q$ en todo nodo (autonom\'ia)};
  \end{scope}
\end{tikzpicture}

\begin{tikzpicture}[
  >={Stealth[length=2mm]},
  depot/.style={rectangle,draw,thick,fill=gray!20,minimum size=7mm,font=\small\bfseries},
  cust/.style={circle,draw,thick,fill=blue!10,minimum size=7mm,font=\small},
  stat/.style={regular polygon,regular polygon sides=3,draw,thick,fill=green!18,
               minimum size=8mm,inner sep=0pt,font=\scriptsize},
  arc/.style={->,very thick,gray!60!black}
]
  % ---------- Ruta (parte superior) ----------
  \node[depot] (d0) at (0,0)   {0};
  \node[cust]  (c1) at (2.2,0.7) {$c_1$};
  \node[cust]  (c2) at (4.4,1.0) {$c_2$};
  \node[stat]  (f)  at (6.6,0.3) {$f$};
  \node[cust]  (c3) at (8.8,0.8) {$c_3$};
  \node[depot] (d1) at (11,0)  {0};

  \draw[arc] (d0) -- (c1);
  \draw[arc] (c1) -- (c2);
  \draw[arc] (c2) -- (f);
  \draw[arc] (f)  -- (c3);
  \draw[arc] (c3) -- (d1);

  % etiquetas entrega (delivery) y recogida (pickup) simultaneas
  \node[font=\scriptsize,blue!50!black,above=0.5mm of c1] {$\delta_1\!\downarrow\ \rho_1\!\uparrow$};
  \node[font=\scriptsize,blue!50!black,above=0.5mm of c2] {$\delta_2\!\downarrow\ \rho_2\!\uparrow$};
  \node[font=\scriptsize,blue!50!black,above=0.5mm of c3] {$\delta_3\!\downarrow\ \rho_3\!\uparrow$};
  \node[font=\scriptsize,green!40!black,below=0.8mm of f] {recarga $q_f$};
  \node[font=\scriptsize,below=0.8mm of d0] {dep\'osito};

  % leyenda compacta
  \node[font=\scriptsize,anchor=west] at (0,2.0)
    {$\downarrow$ entrega \quad $\uparrow$ recogida \quad (simult\'aneas, una sola visita)};

  % ---------- Curva de estado de carga SoC (parte inferior) ----------
  \begin{scope}[yshift=-4.7cm]
    \draw[->] (0,0) -- (11.3,0) node[right,font=\scriptsize] {avance de la ruta};
    \draw[->] (0,0) -- (0,2.3) node[above,font=\scriptsize] {SoC $y$};
    \draw[dashed,gray] (0,2.0) -- (11,2.0) node[right,font=\scriptsize,black] {$Q$};

    % SoC: parte llena en 0, baja con la distancia, salta en la estacion f, vuelve a bajar
    \draw[very thick,green!45!black]
      (0,2.0) -- (2.2,1.45) -- (4.4,0.95) -- (6.6,0.45) % consumo hasta la estacion
      -- (6.6,1.75)                                      % salto por recarga parcial q_f
      -- (8.8,1.15) -- (11,0.55);                        % consumo hasta el deposito

    % marcas verticales de los nodos
    \foreach \x/\lab in {0/0,2.2/c_1,4.4/c_2,6.6/f,8.8/c_3,11/0}
      \draw[gray!40] (\x,0) -- (\x,-0.08) node[below,font=\scriptsize,black] {$\lab$};

    \draw[<->,green!40!black] (6.75,0.45) -- (6.75,1.75)
      node[midway,right,font=\scriptsize] {$q_f$};
    \node[font=\scriptsize,anchor=west] at (0.2,-0.7)
      {restricci\'on: $\,0 \le y \le Q$ en todo nodo (autonom\'ia)};
  \end{scope}
\end{tikzpicture}
